% explainer-source-hash: f853827457ce4135f852201be81a6e16bb0e75ff89457f8e3540348443911a82
% explainer-source-commit: a0adde8872a5b2ad97203deeb662f65644532a78
% explainer-generated: 2026-07-16
% Regenerate with /explain-all (see WIRING.md). Do not edit this stamp by hand:
% it is how the toolchain knows whether this document still matches the code.
\documentclass[11pt,a4paper]{article}
\input{preamble}

\begin{document}

\begin{center}
{\LARGE\bfseries\color{inkblue} urdu-slm\par}
\vspace{0.3cm}
{\large\color{accent} Brief\par}
\vspace{0.5cm}
\begin{minipage}{0.88\textwidth}
\small\itshape
A small Urdu language model trained from scratch: what it is, why it exists, what it
proved, and what is wrong with it. The companion Deep Dive covers every module and
function; this document is the argument, at altitude. A glossary of every technical term
is at the end.
\end{minipage}
\end{center}
\vspace{0.4cm}
\hrule
\vspace{0.6cm}

% =====================================================================
\section{The claim, and whether it survives contact}

Most "train a GPT from scratch" projects use English data and somebody else's tokenizer.
This one picks a low-resource language and refuses both, which forces it to solve the
problems the tutorials skip: script normalization, deduplication of a corpus that is
mostly bot-generated stubs, language identification without a 126MB model, and a
tokenizer that does not shred the language into individual bytes.

The headline claim is that GPT-2's tokenizer needs \textbf{5.23x more tokens} for the same
Urdu text. That is not a rhetorical number; it is the project's whole justification, so it
was checked rather than repeated.

\begin{keyidea}{Verified on this machine, 2026-07-16, to the digit}
A real GPT-2 tokenizer was found in the local Hugging Face cache and
\code{eval/tokenizer\_compare.py} was run against it. Every figure in the README's table
reproduced exactly: an 806,495-character held-out Urdu sample costs \textbf{197,758}
tokens with this project's tokenizer (4.078 chars/token) against GPT-2's
\textbf{1,034,831} (0.779 chars/token). Ratio: \textbf{5.23x}.
\end{keyidea}

\begin{figure}[H]
\centering
\begin{tikzpicture}
\begin{axis}[
  width=12.6cm, height=3.6cm,
  xbar, y=9mm, bar width=6mm,
  xmin=0, xmax=1200000,
  xlabel={tokens spent on the same 806,495 characters of Urdu (lower is better)},
  symbolic y coords={gpt2,ours}, ytick=data,
  yticklabels={GPT-2 (50k English BPE), ours (32k Urdu BPE)},
  nodes near coords, nodes near coords align={horizontal},
  nodes near coords style={font=\footnotesize,text=inkblue},
  every axis plot/.append style={fill=accent!55, draw=inkblue},
  tick label style={font=\small}, label style={font=\small},
  xmajorgrids, grid style={linegrey}, scaled x ticks=false,
  xticklabel style={/pgf/number format/fixed, /pgf/number format/1000 sep={,}},
  enlarge y limits=0.7,
]
\addplot coordinates {(1034831,gpt2) (197758,ours)};
\end{axis}
\end{tikzpicture}
\caption{The measured compression gap, reproduced from the committed tokenizer and
validation data.}
\end{figure}

\subsection{Why the gap is that large}

Two effects compound. UTF-8 gives every ASCII character one byte, but every Urdu letter
lives in the Arabic block and costs \textbf{two}. And GPT-2's 50,257 merges were learned
on English web text, so essentially none of them fire on Urdu byte sequences: the
tokenizer falls back to its base alphabet and emits roughly one token per byte. At 0.779
chars/token it is spending more than one token \emph{per Urdu character}.

Because every cost in the system is counted in tokens, this is not an aesthetic
complaint. The GPU bill is per token. The context window is 1024 tokens, which holds about
4,176 Urdu characters with this tokenizer and about 798 with GPT-2's. And a model whose
tokens are mostly raw bytes must spend its early capacity learning to reassemble
characters before it can learn any Urdu at all.

\begin{honest}{The README's own command does not produce the README's number}
The measurement is documented as \code{python -m eval.tokenizer\_compare --gpt2 <path>},
but \code{--max-chars} defaults to 500,000. Run exactly as written it yields a
504,251-character sample and \textbf{5.14x}. Reproducing 5.23x needs
\code{--max-chars 800000}. Both numbers are honest and the finding is identical either
way; the defect is that the command printed beside the table is not the command that
produced it. One flag fixes it, and it is worth fixing precisely because the number is
genuine.
\end{honest}

% =====================================================================
\section{A note on the script}

\begin{honest}{Why no Urdu appears in this document}
Urdu uses the Perso-Arabic script: right to left, with context-dependent letter shapes.
This PDF is built with \code{pdflatex}, which on this machine has no font for that script
and no right-to-left engine; a multi-byte character inside a code listing is a hard build
failure, not a rendering glitch. So Urdu text here is transliterated, described, or given
as codepoints. The real script is in the repository's README and eval set, and renders
fine in any editor. Nothing is hidden, it simply cannot be drawn here.
\end{honest}

% =====================================================================
\section{The system}

\begin{figure}[H]
\centering
\resizebox{\textwidth}{!}{
\begin{tikzpicture}[node distance=8mm]
  \node[extbox,text width=2.7cm] (raw) {public sources\\ \scriptsize Wikipedia dump,\\ Leipzig news,\\ CC-100 crawl};
  \node[box,right=of raw,text width=2.7cm] (clean) {clean\\ \scriptsize NFC, fold variants,\\ script-ratio filter};
  \node[box,right=of clean,text width=2.7cm] (dedup) {dedup\\ \scriptsize blake2b exact +\\ MinHash LSH near};
  \node[box,right=of dedup,text width=2.6cm] (tok) {tokenize\\ \scriptsize 32k byte-level BPE\\ trained on this corpus};
  \node[store,right=of tok,text width=2.4cm] (bin) {\file{train.bin}\\ \scriptsize flat uint16\\ token stream};
  \node[box,below=10mm of bin,text width=2.9cm] (tr) {\code{train.py}\\ \scriptsize AdamW, cosine LR,\\ bf16, token budget};
  \node[box,left=of tr,text width=2.9cm] (m) {124M transformer\\ \scriptsize RMSNorm, RoPE,\\ SwiGLU, tied emb};
  \node[store,left=of m,text width=2.5cm] (out) {\file{runs/base\_v2/}\\ \scriptsize model.pt +\\ metrics.jsonl};
  \node[box,left=of out,text width=2.6cm] (ev) {\code{eval/}\\ \scriptsize perplexity, script\\ mix, cloze};
  \draw[flow] (raw) -- (clean);
  \draw[flow] (clean) -- (dedup);
  \draw[flow] (dedup) -- (tok);
  \draw[flow] (tok) -- (bin);
  \draw[flow] (bin) -- (tr);
  \draw[flow] (tr) -- (m);
  \draw[flow] (m) -- (out);
  \draw[flow] (out) -- (ev);
  \node[lbl,below=1mm of clean] {\scriptsize 31.9\% removed (v1)};
  \node[lbl,below=1mm of tok] {\scriptsize 4.078 chars/token};
\end{tikzpicture}}
\caption{The whole project. Every stage streams to disk; nothing holds the corpus in
memory. Training memory-maps the token stream.}
\end{figure}

Five components, about 2,800 lines of Python including tests:

\begin{itemize}
\item \textbf{The pipeline} streams raw text through normalization, filtering,
      deduplication and a deterministic train/val split, writing sharded JSONL. Each of
      its three stages writes a marker file so a rerun skips completed work.
\item \textbf{The tokenizer} is a 32k byte-level BPE trained on the cleaned corpus,
      \emph{train split only}, so no validation text shapes its merges.
\item \textbf{The model} is a decoder-only transformer written out in one 186-line file
      with no modeling code imported from \code{transformers}. It is the Llama-family
      design, not a GPT-2 clone: RMSNorm, rotary position embeddings, SwiGLU, no biases,
      tied input/output embeddings, and \code{scaled\_dot\_product\_attention} for
      Flash-kernel causal attention.
\item \textbf{Training} budgets in tokens rather than steps, resumes exactly from a
      one-integer data cursor, runs bf16 on GPU and fp32 on CPU, and appends every metric
      to JSONL.
\item \textbf{Evaluation} covers held-out perplexity, a script-mix diagnostic, a
      hand-written 95-item Urdu cloze harness, and the tokenizer comparison above.
\end{itemize}

\subsection{Three architectural decisions worth defending}

\begin{keyidea}{The vocabulary table dominates a small model, and the project says so}
At a 32k vocabulary and 256 dimensions, the embedding table alone is 8.19M parameters;
the six transformer blocks add only 4.82M. So the \code{tiny} preset is 13.01M, not the
"10M" its config comment still claims, and 63\% of it is vocabulary. The author tuned the
blocks to a defensible 13M and documented the split rather than shrinking the vocabulary
to hit a round number, which would have degraded the very tokenizer the project exists to
demonstrate. All four presets' counts were re-measured here and match the README exactly:
13.01M, 35.66M, 123.69M, 336.38M.
\end{keyidea}

\begin{keyidea}{SwiGLU at 8/3 is why \code{base} needs 14 layers, not 12}
SwiGLU uses three weight matrices where a classic MLP uses two, so at the textbook 4x
width it would cost 50\% more parameters. The 8/3 ratio cancels that exactly
($3 \times 8/3 = 8 = 2 \times 4$). The consequence is that each layer is cheaper than a
GPT-2 layer of the same width, so 12 layers do not reach 124M and the preset uses 14.
Depth, not the memorized layer count, is what hits a target size. This was worked out,
not copied.
\end{keyidea}

\begin{keyidea}{The MinHash rewrite: a real 10x on a real bottleneck}
Deduplication matters here more than usual because Urdu Wikipedia is largely
bot-generated near-identical stubs, which is what produced 2.46M paragraphs from 1.73M
articles. The library's per-shingle \code{MinHash.update()} loops in Python and measured
~250 docs/s, extrapolating to over two hours. The rewrite hashes all of a document's
shingles with xxhash and computes the whole signature in two vectorized numpy operations,
reaching ~2400 docs/s.

The insight that makes it safe: MinHash signatures need only be \emph{internally
consistent}, never to match an external standard. That permits swapping datasketch's sha1
for xxhash and writing \code{hashvalues} directly, provided every document shares one
scheme and one set of permutations, which reusing the seed object's own coefficients
guarantees. The LSH banding on top is unchanged.
\end{keyidea}

% =====================================================================
\section{What it actually produced}

Three real runs. Everything in this section was read out of the committed
\file{metrics.jsonl} logs, not out of the README.

\begin{table}[H]
\centering
\small
\begin{tabular}{lrrr}
\toprule
 & \textbf{tiny (CPU)} & \textbf{base v1} & \textbf{base v2} \\
\midrule
Parameters        & 13.01M & 123.69M & 123.69M \\
Steps             & 500 & 858 & 4,768 \\
Tokens seen       & 4.10M & 449.8M & 2,499,805,184 \\
Unique tokens     & slice of v1 & 111.7M & 862.4M \\
Tokens/parameter  & n/a & 0.9 & 20.2 \\
Throughput        & ~1,600 tok/s & ~99,470 tok/s & ~105,000 tok/s \\
Train loss        & 9.74 to 5.42 & 9.59 to 3.51 & 9.63 to 2.74 \\
Held-out perplexity & 214.11 (re-measured) & 38.02 & \textbf{28.93} \\
Cost              & free & ~\$1.90 & ~\$13.06 \\
\bottomrule
\end{tabular}
\caption{The three runs. Perplexity for the two base models is README-attributed (see the
verification note below); everything else was read from the logs.}
\end{table}

\begin{figure}[H]
\centering
\begin{tikzpicture}
\begin{axis}[
  width=13.4cm, height=5.8cm,
  xlabel={step (524,288 tokens each)}, ylabel={train loss (nats)},
  xmin=0, xmax=4850, ymin=2.3, ymax=10,
  tick label style={font=\small}, label style={font=\small},
  xmajorgrids, ymajorgrids, grid style={linegrey},
  legend style={font=\footnotesize, at={(0.98,0.95)}, anchor=north east, draw=linegrey},
]
\addplot[draw=accent, thick, mark=none] coordinates {
(20,9.629) (140,6.831) (260,5.678) (380,5.027) (500,4.627) (620,4.235) (740,3.764)
(860,3.469) (980,3.174) (1100,3.200) (1220,3.154) (1340,3.052) (1460,3.000)
(1580,2.825) (1700,2.933) (1820,2.803) (1940,2.791) (2060,2.808) (2180,2.731)
(2300,2.725) (2420,2.686) (2540,2.675) (2660,2.718) (2780,2.704) (2900,2.625)
(3020,2.682) (3140,2.664) (3260,2.584) (3380,2.665) (3500,2.688) (3620,2.572)
(3740,2.583) (3860,2.605) (3980,2.584) (4100,2.496) (4220,2.612) (4340,2.550)
(4460,2.469) (4580,2.729) (4700,2.773) (4760,2.741) };
\draw[warnred, thick, dashed] (axis cs:4570,2.3) -- (axis cs:4570,10);
\node[lbl, anchor=south east, text=warnred] at (axis cs:4550,6.4)
  {anneal starts: switches to\\ Wikipedia-only data};
\legend{train loss, base v2}
\end{axis}
\end{tikzpicture}
\caption{The v2 run, from the committed log. The step at the anneal boundary is real
data. Its position corroborates the README's account: the main phase was ~2.4B tokens,
which at 524,288 tokens/step lands at step ~4,578, and the log's jump is between 4,560
and 4,580.}
\end{figure}

\subsection{The most important number is one the project refuses to claim}

v2 gave the model 7.7x more unique tokens than v1 and moved the token/parameter ratio
from 0.9 to 20.2, right at the Chinchilla-optimal ~20. Perplexity fell from 38.02 to
28.93, a genuine ~24\% reduction. Cloze accuracy went from 34.9\% to 39.5\%.

And the README declines to call that a cloze improvement, because it is a 2-item
difference on a 43-item set whose standard error is roughly $\pm$7 points. It sits inside
the noise band, and the README says so instead of presenting "+4.6 points" as a result.

\begin{keyidea}{The diagnosis, and why it is the best thinking in the project}
The explanation given is correct and well-sourced: a 124M-parameter model has a real
ceiling on how many discrete facts it can store, and that capacity scales with
\emph{parameter count}, not with training-text volume (Roberts et al., 2020). So 7.7x
more text made the model more fluent, which is exactly what perplexity measures, without
making it know more, which is what cloze measures. The corpus expansion did not fail; it
did precisely what more text does.

That diagnosis then drives the plan. The \code{medium} preset (336.38M, 2.7x the
parameters) exists because parameters are the lever that moves fact storage. The Wikidata
pipeline exists to supply dense declarative facts rather than diffuse prose. And the eval
set was widened from 45 to 95 items \emph{first}, cutting the noise band from $\pm$7 to
$\pm$5.2, because the old set could not resolve the difference the next run is meant to
produce.

Noticing that your headline metric is too noisy to support your next claim, and fixing
the metric before running the experiment, is the most mature decision in this repository.
\end{keyidea}

The v2 model writes essentially perfect Urdu script (99.75\% of generated characters are
Arabic script over 50 generations), which says the corpus filtering worked. What it does
not have is reliable prompt-following: sample generations are fluent but drift off topic,
and the README says so plainly rather than cherry-picking. Fixing that would need a
larger model, fact-structured data, or an instruction-tuning pass, none of which this
project attempts.

% =====================================================================
\section{What is wrong with it}

The Deep Dive lists thirty findings. These are the ones that would cost something to be
surprised by. The pattern across almost all of them is the same, and it is an unusual one:
\textbf{the code is in better shape than the prose around it}.

\begin{honest}{Upsampling repeats documents back-to-back inside one context window}
\code{--upsample wikipedia:3} is implemented as \code{buf.extend(ids * reps)}, which is
Python list repetition and lays the three copies down \emph{contiguously} in the token
stream. A typical kept paragraph is well under 100 tokens, so all three copies fit inside
a single 1024-token training block. The model therefore repeatedly sees a context whose
correct continuation is a verbatim copy of text ~70 tokens earlier, which rewards copying
rather than modeling. That is not what upsampling is meant to do; the intent is to make
Wikipedia's \emph{distribution} count for more. It shipped in the run that produced the
headline model. Whether it measurably hurt cannot be established without a retrain. The
fix is to repeat at the document-list level. This is the most substantive code finding.
\end{honest}

\begin{honest}{The planned Wikidata data would contaminate the eval set the run exists to improve}
\code{fetch\_wikidata.py}'s docstring and \file{docs/DATA\_LICENSES.md} both claim the
fact triples use "different entities than the eval set's own items". Nothing enforces
this, and it is false: the \code{capital} query selects the capital of \emph{every
country}, so it returns Pakistan/Islamabad, and cloze item 1 glosses as "The capital of
Pakistan is Islamabad." Item 21 is covered the same way by the official-language query.

Two of 95 items, so the contamination is small, and Urdu Wikipedia already contains that
fact in prose, so the eval was never clean in the strict sense. What is specifically
wrong is a distinctness claim asserted twice and enforced nowhere, landing on the one
metric the \code{medium} run is designed to move. Fixable in about thirty lines, before
it matters.
\end{honest}

\begin{honest}{The shipped config cannot reproduce the shipped model}
\file{configs/base.yaml} says \code{max\_tokens: 450000000} and its comments describe the
superseded 111.7M-token v1 corpus. But \code{runs/base\_v2} trained on 2,499,805,184
tokens, confirmed exactly against its own log ($4{,}768 \times 524{,}288$), plus an anneal
phase on a different data directory. Both were done with CLI overrides that exist in
\code{train.py}, but \textbf{the actual commands are recorded nowhere}, and
\file{docs/GPU\_RUN.md} still documents only \code{python train.py --config
configs/base.yaml}, which would train a different model on a different budget. Someone
following the documentation cannot reproduce the model the repository is about.
\end{honest}

\begin{honest}{One number does not reproduce: the tiny run's perplexity}
The README reports the tiny model's held-out perplexity as 248.8. The committed script,
on the committed checkpoint and the committed validation set, gives \textbf{214.11} at
the default coverage and \textbf{177.01} over the entire split. The script is fully
deterministic, so this is not noise, and no batch budget produces 248.8.

The likely cause is visible in the timestamps: \file{runs/tiny/model.pt} was written at
20:21, but \file{data/val.bin} was regenerated at 21:46, over an hour later. The figure
was measured against a validation set that no longer exists in the repository. Nothing was
inflated (the real number is better than the claim), but in a project whose entire pitch
is honest reporting, the one figure that does not reproduce is worth re-running.
\end{honest}

\begin{honest}{\file{docs/GPU\_RUN.md} is a pre-run plan whose predictions were falsified}
It estimates throughput at "a conservative 30k tok/s" and states that "the only measured
throughput in this repo is the CPU proof run". Two GPU runs have since measured 99,470
and ~105,000 tok/s: the estimate was low by a factor of ~3.3, and every cost projection
in the document inherits the error. Its token-budget reasoning is v1-era throughout, and
its own advice ("add more data before spending GPU money") was taken and completed.

Its post-run half is excellent, and that is why this is worth fixing rather than deleting:
the RunPod SSH-key injection trap, the warning that installing the pinned
\code{torch==2.12.1} silently replaces a pod's CUDA build with a CPU-only wheel, the
\code{tmux} detach pattern, and the note that SSH cannot stop the billing meter so you
must Terminate rather than Stop. That is hard-won operational knowledge. The document is
two documents stapled together; splitting them would turn the weakest artifact into one of
the strongest.
\end{honest}

\subsection{Shorter}

\begin{itemize}
\item \textbf{\code{generate} has no KV cache}, so sampling is quadratic. Fine for a few
      80-token diagnostics, disqualifying for serving. It is also why the
      grouped-query-attention branch has nothing to optimize: GQA shrinks a KV cache, and
      there is none. No preset enables GQA anyway, and its \code{repeat\_interleave}
      implementation materializes the expanded tensors, so it saves no memory. The
      README's word "capable" is exactly right.
\item \textbf{The DDP path has never been run.} Written carefully, including the
      gradient-sync suppression that cuts inter-GPU traffic 16x, but unproven.
\item \textbf{The LSH index is entirely in RAM}, the one place the streaming discipline
      breaks. The README names this itself.
\item \textbf{v1 never ran a single validation check}, because \code{eval\_interval}
      (1000) exceeded the run's total steps (858).
\item \textbf{All 95 cloze items have the answer at index 0.} Disclosed, and the harness
      scores log-probabilities without looking at position, so there is no leak. But it
      makes a whole class of bug undetectable: a position-bias regression would score
      100\% and look like a triumph.
\item \textbf{Docstrings promise slightly more than the code delivers.}
      \code{train.py} claims resume restores RNG (the checkpoint contains no RNG state;
      near-harmless, since the data order comes from the restored cursor and dropout is 0).
      \code{cloze.py} claims "total log-probability" but computes a length-normalized
      mean, which is the better choice, undocumented.
\item \textbf{Stale docs}: \file{evaldata/README.md} still says 45 items,
      \file{tiny.yaml} still says "~10M params", \file{base.yaml} still says "24GB GPU",
      and \file{data/corpus\_stats.json} describes v1 only, with no committed stats for
      the v2 corpus the headline model trained on.
\item \textbf{\code{tqdm} is declared in \file{requirements.txt} and never imported.}
\end{itemize}

\begin{honest}{What could not be verified here}
The base models' downstream metrics (perplexity 38.02 and 28.93, script mix, cloze
accuracy), the v2 corpus table, the GPU costs, and the Wikidata fetch counts are
README-attributed. \file{data/train.bin} is gitignored (1.7GB, over GitHub's limit) and
the v2 validation set is not committed, so those evaluations cannot be re-run without
re-downloading ~1.3GB of sources and hours of pipeline. What \emph{was} verified: the
tokenizer's headline number to the digit, all 13 tests passing, all four presets'
parameter counts, the v1 corpus funnel, all three runs' losses and throughputs, the v2
token count exactly, the anneal boundary's position, the v2 validation curve, and the
eval set's shape.
\end{honest}

% =====================================================================
\section{The verdict}

\begin{keyidea}{What this project demonstrates}
The claim it is built on is real and reproduces exactly. The engineering underneath is
sound and occasionally sharp: the MinHash vectorization is a genuine 10x with an
articulable safety argument; the SwiGLU sizing was reasoned out rather than copied;
training budgets in tokens so the scaling reasoning is expressible in the config; the
causal-masking test tests the property by experiment and includes the negative assertion
that stops it being a tautology; validation is never upsampled and the tokenizer never
sees it.

The Urdu-specific decisions require actually knowing something about the orthography
rather than about machine learning: keeping the zero-width non-joiner (it carries
meaning) while stripping its invisible neighbours, and refusing to strip harakat because
they matter in the minority of text that has them.

And the reporting is unusually honest, including where it is unflattering: the "10M"
model is 13M and the README says why, the cloze gain is inside the noise band and the
README refuses to claim it, and the eval set was widened to fix the noise \emph{before}
running the experiment that needs it.

The weaknesses are almost entirely documentation drifting behind code: a GPU plan never
updated after the run falsified it, configs that no longer describe the model that
shipped, docstrings overreaching slightly, and a perplexity measured against a validation
set that was later regenerated. Two real code defects (contiguous upsampling, and the
unenforced Wikidata distinctness claim) are worth fixing before the next run. That is the
opposite of the usual failure mode, and a much easier problem to have.
\end{keyidea}

% =====================================================================
\section{Glossary}

\begin{description}[leftmargin=0pt,style=nextline,itemsep=3pt]
\item[Annealing] Finishing training on a small, high-quality data slice after a large
  mixed corpus, so the final steps bias the model toward the better distribution. v2 ended
  with a ~100M-token Wikipedia-only phase; the switch is visible as a loss step in the log.
\item[Attention] The mechanism letting each position look at others. Each position emits a
  query, a key and a value; queries are compared against all keys, and the resulting
  weights average the values.
\item[bfloat16 / AMP] A 16-bit float keeping float32's exponent range but fewer mantissa
  bits. It cannot overflow where fp32 would not, so it needs no loss scaling. Automatic
  Mixed Precision runs most operations in it while keeping master weights in fp32.
\item[BPE (Byte-Pair Encoding)] Tokenizer training: start from single bytes and repeatedly
  merge the most frequent adjacent pair into a new vocabulary entry until reaching the
  target size (32,000 here). \emph{Byte-level} BPE seeds the vocabulary with all 256 bytes,
  so no input can ever be out-of-vocabulary.
\item[Causal masking] Restricting each position to attend only to itself and earlier
  positions. Without it the model could read the answer it is asked to predict.
\item[Chinchilla-optimal] The rule of thumb that compute-optimal training uses ~20 tokens
  per parameter. v1 was at 0.9 (badly data-limited); v2 reached 20.2.
\item[Cloze] A fill-in-the-blank test. The model scores each candidate completion rather
  than generating; the best-scoring one is its answer.
\item[Context window] The maximum tokens the model attends over at once: 1024 here.
  Attention cost grows with its square.
\item[Cross-entropy loss] The training objective, in nats: $-\ln p$ for the true next
  token. Lower means less surprised.
\item[DDP] DistributedDataParallel: each GPU holds a full model copy and a different batch
  slice; gradients are averaged after each backward pass.
\item[Deduplication, exact and near] Exact uses a blake2b hash. Near uses MinHash+LSH to
  catch documents that differ slightly, which exact hashing misses entirely.
\item[Embedding] The lookup table turning a token id into a vector. At 32k vocabulary it
  dominates small models: 63\% of \code{tiny}.
\item[FlashAttention] An exact attention implementation that never writes the full
  position-by-position matrix to GPU memory, tiling it through fast on-chip memory instead.
\item[Gradient accumulation] Running several micro-batches that each add into the same
  gradient buffers before one optimizer step, giving a large effective batch without the
  VRAM.
\item[GQA (grouped-query attention)] Sharing keys and values across groups of heads to
  shrink the generation-time KV cache. Implemented here but enabled by no preset.
\item[Jaccard similarity] Intersection over union of two sets. Here, of two documents'
  shingle sets. The near-dedup threshold is 0.8.
\item[JSONL] One complete JSON object per line; readable and appendable one record at a
  time.
\item[KV cache] Storing past keys and values during generation so each new token costs
  $O(1)$ rather than a full re-run. This project has none, so sampling is quadratic.
\item[LSH] Locality-Sensitive Hashing: bucket signatures so only plausibly-similar
  documents are ever compared, avoiding a quadratic scan.
\item[Memory-mapping] Presenting a file as if it were an in-memory array, loading pages on
  demand. How a 1.7GB corpus trains without a 1.7GB allocation.
\item[MinHash] A fixed-size signature (64 numbers here) whose agreement rate estimates
  Jaccard similarity, computed without comparing the sets.
\item[Nats] The unit of loss when using natural logarithms. Perplexity is $e^{\mathrm{nats}}$.
\item[Perplexity] $e^{\mathrm{loss}}$. Roughly "how many equally-likely options did the
  model feel it was choosing between". v2 scores 28.93.
\item[Pre-norm] Normalizing a sublayer's input rather than its output, leaving the residual
  path clean. Markedly more stable than post-norm.
\item[Residual connection] \code{x = x + sublayer(norm(x))}: the sublayer computes a
  correction rather than a replacement, giving gradients a clean path to early layers.
\item[RMSNorm] Normalization dividing by the root mean square with one learned scale,
  dropping LayerNorm's mean-subtraction and shift. Computed in fp32 here for stability.
\item[RoPE] Rotary position embeddings: rotate queries and keys (never values) by an angle
  proportional to position, so attention scores depend only on relative distance. Zero
  parameters.
\item[Shingle] A short overlapping substring (5 characters here). Documents sharing most
  shingles are near-duplicates.
\item[SwiGLU] A gated MLP: the input is projected twice, one path passes through SiLU and
  multiplies the other. At an 8/3 ratio it matches a 4x GELU block's parameter count.
\item[Token] The unit a model reads and writes: a frequently occurring chunk of characters,
  not a letter and usually not a word.
\item[Weight tying] Using one matrix for both the input embedding and the output
  projection. Saves 24.6M parameters in \code{base} and usually improves quality.
\item[ZWNJ] Zero-width non-joiner (U+200C): invisible, but carries real orthographic
  meaning in Urdu. Deliberately kept where its invisible neighbours are stripped.
\end{description}

\end{document}
